\documentclass{article}
\usepackage{amsmath}
\usepackage{enumitem}
\title{Quantitative Methods for Business \\[1em] \large Extra Assignment}
\date{\today}
\author{Daniel Sánchez Pazmiño}

\begin{document}
\maketitle

\begin{enumerate}
\item Write down the gradient and the Hessian for each of the following functions:
\end{enumerate}

\begin{itemize}
    \item[(a)] $f(x, y) = x^4 + x^2 - 6xy + 3y^2$
    \item[(b)] $f(x, y) = xy^2 + x^3y - xy$
    \item[(c)] $f(x, y, z) = x^2 + 6xy + y^2 - 3yz + 4z^2 - 10x - 5y - 21z$
\end{itemize}

\textit{Note}: The Hessian collects the second partial derivatives, as follows:

\[
H = \begin{bmatrix}
    \dfrac{\partial^2 f}{\partial x^2} & \dfrac{\partial^2 f}{\partial x \partial y} \\
    \dfrac{\partial^2 f}{\partial y \partial x} & \dfrac{\partial^2 f}{\partial y^2}
\end{bmatrix}
\]

\begin{enumerate}[resume]
    \item Consider the function $f(x) = - \dfrac{x^3}{3} + x^2 - x +4$.
    \begin{enumerate}
        \item Draw the graph of $f(x)$.
        \item Find the critical points of $f(x)$.
        \item Determine the concavity of $f(x)$.
        \item Find any local maxima or minima of $f(x)$, if they exist.
        \item If I were to restrict the domain of $f(x)$ to [3, 3] what would be the maximum and minimum values of $f(x)$? Are these values local maxima or minima?
    \end{enumerate}
\end{enumerate}

\begin{enumerate}[resume]
    \item Consider demand for hospital beds in a city, given by $Q = 100 - 2P$, where $Q$ is the quantity of hospital beds demanded and $P$ is the price of a hospital bed. 
    \begin{enumerate}
        \item When COVID-19 hit the city, the demand for hospital beds increased and shifted right to $Q = 200 - 2P$. Draw the graph of the demand curve before and after COVID-19.
        \item If the supply of hospital beds is given by $Q = 50 + 3P$, find the equilibrium price and quantity of hospital beds before and after COVID-19. 
        \item If the city government wants to increase the quantity of hospital beds to 300, what should the price be?
        \item A reduced form estimation of demand for hospital beds in the city is given by $Q = 100 - 2P + 0.5I$, where $I$ is the income of the city. Calculate the partial price elasticity of demand for hospital beds in the city at equilibrium, with an Income $I$ of 100. Interpret the result.
        \item Assume the government intervenes and increases the price by 20\%. Calculate the percentage change in quantity demanded of hospital beds. Interpret the result. 
    \end{enumerate}
\end{enumerate}

\begin{enumerate}[resume]
    \item Consider the sum of squared errors, a statistical measure that will be covered later in the program. It is expressed as $$SSE = \sum_{i=1}^{n} (y_i - \hat{y}_i)^2$$

    where $y_i$ is the actual value of the dependent variable, $\hat{y}_i$ is the predicted value of the dependent variable, and $n$ is the number of observations.
\begin{enumerate}
    \item Using properties of sums, show that the sum of squared errors can be expressed as $$SSE = \sum_{i=1}^{n} y_i^2 - 2\sum_{i=1}^{n} y_i\hat{y}_i + \sum_{i=1}^{n} \hat{y}_i^2$$
\end{enumerate}
\end{enumerate}


\end{document}
